\section{Exponentes enteros negativos}

La regla del cociente nos permite simplificar, por ejemplo,
\[
\frac{2^5}{2^3}=2^{5-3}=2^2.
\]
Pero ¿qué ocurre si el exponente del numerador es menor que el del denominador? Considere
\[
\frac{2^3}{2^5}.
\]

Desarrollemos primero las potencias:
\[
\frac{2^3}{2^5}
=
\frac{2\cdot2\cdot2}
{(2\cdot2\cdot2)(2\cdot2)}.
\]

Los tres factores del numerador también aparecen como un factor del denominador. Podemos agruparlos:
\[
\frac{2^3}{2^5}
=
\frac{2^3}{2^3\cdot2^2}.
\]

Como \(2^3\ne0\), el cociente de \(2^3\) consigo mismo es el elemento neutro multiplicativo:
\[
\frac{2^3}{2^3}=1.
\]
Por lo tanto,
\[
\frac{2^3}{2^3\cdot2^2}
=
\frac{1}{2^2}.
\]

Observe que ningún factor simplemente ``desapareció''. Los factores comunes produjeron un \(1\), y como multiplicar por \(1\) no cambia el valor de una expresión, solemos dejar de escribirlo.

Así obtenemos
\[
\frac{2^3}{2^5}
=
\frac{1}{2^2}.
\]

Ahora bien, si queremos conservar la misma regla de restar exponentes, también deberíamos tener
\[
\frac{2^3}{2^5}=2^{3-5}=2^{-2}.
\]
Las dos expresiones deben representar entonces el mismo número:
\[
2^{-2}=\frac1{2^2}=\frac14.
\]
